Filaments carry much of the cross-field transport in the tokamak scrape-off layer, and models predict how their radial velocity scales with their size. We measure both on MAST from the 100~kHz, $48 \times 256$ pixel strip recordings of the midplane fast camera in the public FAIR-MAST archive. An automatic pipeline built from standard computer-vision operations (running-minimum background, robust thresholding, connected components, Hungarian linking, least-squares tracks, and optical flow) produces \nTracksSel{} filament tracks in \nShotsStats{} discharges across four campaigns. The archive distributes no spatial calibration for these recordings; we recover the pixel scale for each camera configuration by regressing the image column of the plasma limb on the EFIT outboard LCFS radius across discharges, with bootstrap intervals, and obtain \scaleOwnMin{}--\scaleOwnMax{}~mm per pixel. The median radial velocity is \vMedAll{}~m/s and the median radial width is \dMedAll{}~cm. After correction for exposure blur, the velocity--width relation has a log--log slope of \vdSlopePooledBlur{} (95\% interval \vdSlopePooledBlurLo{} to \vdSlopePooledBlurHi), so the velocity is nearly independent of size over the measured range. In two-region units the filaments span $\hat a = \aHatMin$--$\aHatMax$, straddling the model's velocity maximum, which accounts for the flat dependence; but their normalized velocity is $\hat v \approx \vHatMed$, far below the model scale. Code and derived data are released.
